\documentclass{article}
\usepackage{amsmath,amssymb}
\usepackage{xurl}
\usepackage[hidelinks]{hyperref}
\setlength{\emergencystretch}{2em}
\interfootnotelinepenalty=10000
\widowpenalty=10000
\clubpenalty=10000
% Shared commands for notes in docs/paper-gaps/.

\DeclareUnicodeCharacter{2080}{\ifmmode{}_{0}\else\textsubscript{0}\fi}
\DeclareUnicodeCharacter{2081}{\ifmmode{}_{1}\else\textsubscript{1}\fi}
\DeclareUnicodeCharacter{2082}{\ifmmode{}_{2}\else\textsubscript{2}\fi}
\DeclareUnicodeCharacter{2099}{\ifmmode{}_{n}\else\textsubscript{n}\fi}

\newcommand{\Lean}{\textsc{Lean}}
\ExplSyntaxOn
\NewDocumentCommand{\leanid}{m}{
  \ifmmode
    \textsf{\detokenize{#1}}
  \else
    \group_begin:
    \urlstyle{sf}
    \tl_set:Nn \l_tmpa_tl {#1}
    \tl_replace_all:Nnn \l_tmpa_tl {\_} {_}
    \exp_args:NV \path \l_tmpa_tl
    \group_end:
  \fi
}
\ExplSyntaxOff
\providecommand{\tr}{\operatorname{tr}}
\newcommand{\tnleanIssueBase}{https://github.com/LionSR/TNLean/issues/}
\newcommand{\tnleanPrBase}{https://github.com/LionSR/TNLean/pull/}
\newcommand{\tnleanDocBase}{https://sirui-lu.com/TNLean/docs/}
\newcommand{\ghissue}[1]{\href{\tnleanIssueBase#1}{GitHub issue \##1}}
\newcommand{\ghpr}[1]{\href{\tnleanPrBase#1}{PR \##1}}
\newcommand{\doclink}[1]{\href{\tnleanDocBase#1.html}{\path{#1}}}

% Dirac notation and the identity operator, as in blueprint/src/macros/common.tex.
% \providecommand keeps note-local definitions authoritative.
\usepackage{dsfont}
\providecommand{\ket}[1]{|#1\rangle}
\providecommand{\bra}[1]{\langle #1|}
\providecommand{\braket}[2]{\langle #1 | #2 \rangle}
\providecommand{\ketbra}[2]{|#1\rangle\!\langle #2|}
\providecommand{\Id}{\mathds{1}}
\providecommand{\one}{\mathds{1}}
\providecommand{\C}{\mathbb{C}}
\providecommand{\MN}[1]{M_{#1}(\C)}
\providecommand{\MD}{M_D(\C)}

% Verdict marker: \gapnote{<kind>}{<status>}. Kind is the policy
% classification (clarification, local-correction, scope-restriction,
% unfaithful, false-source, open-gap); status is open, wip (elimination
% underway), resolved, or historical. Machine-read by the site index;
% prints a verdict line.
\newcommand{\gapnote}[2]{\par\noindent\textbf{Verdict:} #1 (#2).\par}

\title{Polynomial Approximation Accuracy: Equal Blocks and Positive Correlation Length}
\author{}
\date{2026-10-03}
\begin{document}
\maketitle
\gapnote{scope-restriction}{open}

\section{The source statement}
After Lemma~1, \cite[p.~4]{Malz2024Preparation} states that, for any
$\eta>0$, blocking
\begin{equation}
    q=\left\lceil 2\xi(1+\eta)\log N\right\rceil
    \label{eq:mswc24_polynomial_block_length}
\end{equation}
sites gives approximation error $O(N^{-\eta})$. Here the target is the
normalized periodic state of a normal tensor, the error is
$\varepsilon(\psi,\phi)=1-|\langle\psi,\phi\rangle|$, and
$\xi=-1/\log|\lambda_2|$ is the correlation length. The construction uses
the gauge of eq.~(5), with identity left fixed matrix and positive definite
right fixed matrix $\rho$ satisfying $\operatorname{tr}\rho=1$.

\section{The statement proved for equal blocks}
The result in \doclink{TNLean/MPS/Preparation/PolynomialAccuracy} concerns
finite positive correlation length $\xi>0$ and equal blocks. There is
$C>0$, depending on the tensor and its fixed-point data, such that for
every $\eta>0$ and all positive integers $q,M$ with
$N=Mq\geq2$ and (\ref{eq:mswc24_polynomial_block_length}),
\begin{equation}
    1-\left|\left\langle\widetilde\phi_N,\phi_N\right\rangle\right|
    \leq C N^{-\eta}.
    \label{eq:mswc24_polynomial_uniform_error}
\end{equation}
The vector $\widetilde\phi_N$ is the normalized periodic state obtained
by replacing the positive part of the $q$-site blocked tensor by its
fixed-point limit, as in eqs.~(9)--(10) of the source. The constant $C$ is
independent of $\eta,q,M$.

The proof uses the stronger second-order estimate already established in
\doclink{TNLean/MPS/Preparation/ApproximationError}:
\begin{equation}
    \varepsilon(\widetilde\phi_N,\phi_N)
    \leq C_\gamma M e^{-2\gamma q/\xi},\qquad 0<\gamma<1.
    \label{eq:mswc24_polynomial_second_order}
\end{equation}
Fixing $\gamma=1/4$ makes the constant independent of $\eta$. Since
$q\geq2\xi(1+\eta)\log N$ and $q\geq1$, one has $M\leq N$. The bound
(\ref{eq:mswc24_polynomial_uniform_error}) follows from
\begin{equation*}
    M e^{-q/(2\xi)}
    \leq N e^{-(1+\eta)\log N}=N^{-\eta}.
\end{equation*}
The printed Lemma~1 gives only $O(Me^{-\gamma q/\xi})$ for
$\gamma<1/2$. Substituting (\ref{eq:mswc24_polynomial_block_length}) into
that estimate gives $O(N^{1-2\gamma(1+\eta)}/q)$, whose power of $N$
has exponent strictly greater than $-\eta$. The factor $1/q$ decays only
logarithmically for fixed $\eta$. Thus that substitution alone does not
establish the exact coefficient in
(\ref{eq:mswc24_polynomial_block_length}); the proved second-order
estimate (\ref{eq:mswc24_polynomial_second_order}) supplies the required
justification.

\section{The restriction and its elimination}
The equality $N=Mq$ restricts
(\ref{eq:mswc24_polynomial_uniform_error}) to chain lengths divisible by
the prescribed block length. It does not assert that such a decomposition
exists for every $N$ and $\eta$. For example, for fixed $\eta$ and
sufficiently large prime $N$, the prescribed $q$ lies strictly between
$1$ and $N$ and does not divide $N$.

The divisibility restriction has been removed for the logarithmic-depth
preparation theorem \cite{gap:mswc24_depth_upper_bound_divisible_length},
but that result does not by itself prove the polynomial error bound with
the exact block length (\ref{eq:mswc24_polynomial_block_length}). Its
unequal-block error estimate has the slower rate of the printed Lemma~1.
To remove the present restriction, extend the second-order estimate to
blocks of lengths $q,\ldots,q,q'$ with $q\leq q'<2q$. The same exponential
comparison then applies when $N\geq q$. For $N<q$, use exact preparation
of the short chain, as in the logarithmic-depth theorem. Zero correlation
length requires a separate fixed-point argument and is not covered by
the hypothesis $\xi>0$ here.

\bibliographystyle{alpha}
\bibliography{references}
\end{document}
